\clearpage
\onecolumn
\supplementarysection
% Supplementary material.  The figures and the files in generated/ are made by paper/build_paper.py from the results
% of the pipeline (results/flux_model); rerun it after a new run.  \SIdir is the folder of this file relative to the
% main .tex file (define it before \input{...supplementary} if the folder is elsewhere).
\providecommand{\SIdir}{supplementary/}
\input{\SIdir generated/numbers}
\subsection{Kinetic model expressions}\label{sec: kinetic expressions}



% ---------------------------------------------------------------
\textbf{Phosphotransferase system (PTS)}

\begin{align}
    v_1 =
    \frac{
        v^{\max}_{PTS}\,[glc]\,\dfrac{[pep]}{[pyr]}
    }{
        K_{a1} + K_{a2}\dfrac{[pep]}{[pyr]}
        + K_{a3}[glc]
        + [glc]\,\dfrac{[pep]}{[pyr]}
    }
    \cdot
    \frac{1}{1 + \left(\dfrac{[g6p]}{K_{g6p}}\right)^{n_{g6p}}}
    \label{eqn: PTS}
\end{align}
Empirical formulation from \textcite{Kadir2010ModelingsimulationmainmetabolismEscherichiacoliitsseveralsinglegeneknockoutmutantsexperimentalverification}; thermodynamic driving force not applicable.

% ---------------------------------------------------------------
\textbf{Glucose 6-phosphate isomerase (PGI)} \quad EC: 5.3.1.9 \quad $g6p \rightleftharpoons f6p$

\begin{equation}
    \kappa_2 = \frac{[g6p]/K_{S,g6p}^2}%
                    {1 + [g6p]/K_{S,g6p}^2 + [f6p]/K_{P,f6p}^2}
    \qquad
    \gamma_2 = 1 - \frac{[f6p]}{[g6p]\,K'_{eq,2}}
\end{equation}
\begin{equation}
    v_2 = Pgi \cdot k^+_{cat,2} \cdot \kappa_2 \cdot \gamma_2
    \label{eqn: pgi}
\end{equation}

% ---------------------------------------------------------------
\textbf{Phosphofructokinase (PFK)} \quad EC: 2.7.1.11 \quad $f6p + atp \rightleftharpoons fbp + adp$

\begin{equation}
    \kappa_3 = \frac{([f6p]/K_{S,f6p}^3)\,([atp]/K_{S,atp}^3)}%
                    {1 + ([f6p]/K_{S,f6p}^3)([atp]/K_{S,atp}^3) + ([fbp]/K_{P,fbp}^3)([adp]/K_{P,adp}^3)}
    \qquad
    \gamma_3 = 1 - \frac{[fbp]\,[adp]}{[f6p]\,[atp]\,K'_{eq,3}}
\end{equation}
\begin{equation}
    v_3 = PfkB \cdot k^+_{cat,3} \cdot \kappa_3 \cdot \gamma_3
    \label{eqn: pfk}
\end{equation}

% ---------------------------------------------------------------
\textbf{Fructose-bisphosphate aldolase (FBA)} \quad EC: 4.1.2.13 \quad $fbp \rightleftharpoons g3p + dhap$

\begin{equation}
    \kappa_4 = \frac{[fbp]/K_{S,fbp}^4}%
                    {1 + [fbp]/K_{S,fbp}^4 + ([g3p]/K_{P,g3p}^4)([dhap]/K_{P,dhap}^4)}
    \qquad
    \gamma_4 = 1 - \frac{[g3p]\,[dhap]}{[fbp]\,K'_{eq,4}}
\end{equation}
\begin{equation}
    v_4 = FbaA \cdot k^+_{cat,4} \cdot \kappa_4 \cdot \gamma_4
    \label{eqn: fba}
\end{equation}

% ---------------------------------------------------------------
\textbf{Triosephosphate isomerase (TPI)} \quad EC: 5.3.1.1 \quad $dhap \rightleftharpoons g3p$

\begin{equation}
    \kappa_5 = \frac{[dhap]/K_{S,dhap}^5}%
                    {1 + [dhap]/K_{S,dhap}^5 + [g3p]/K_{P,g3p}^5}
    \qquad
    \gamma_5 = 1 - \frac{[g3p]}{[dhap]\,K'_{eq,5}}
\end{equation}
\begin{equation}
    v_5 = TpiA \cdot k^+_{cat,5} \cdot \kappa_5 \cdot \gamma_5
    \label{eqn: tpi}
\end{equation}

% ---------------------------------------------------------------
\textbf{Glyceraldehyde 3-phosphate dehydrogenase (GAPDH)} \quad EC: 1.2.1.12 \quad $g3p + pi + nad \rightleftharpoons pgp + nadh$

\begin{equation}
    \kappa_6 = \frac{([g3p]/K_{S,g3p}^6)([pi]/K_{S,pi}^6)([nad]/K_{S,nad}^6)}%
                    {1 + ([g3p]/K_{S,g3p}^6)([pi]/K_{S,pi}^6)([nad]/K_{S,nad}^6) + ([pgp]/K_{P,pgp}^6)([nadh]/K_{P,nadh}^6)}
\end{equation}
\begin{equation}
    \gamma_6 = 1 - \frac{[pgp]\,[nadh]}{[g3p]\,[pi]\,[nad]\,K'_{eq,6}}
\end{equation}
\begin{equation}
    v_6 = GapA \cdot k^+_{cat,6} \cdot \kappa_6 \cdot \gamma_6
    \label{eqn: gapdh}
\end{equation}

% ---------------------------------------------------------------
\textbf{Phosphoglycerate kinase (PGK)} \quad EC: 2.7.2.3 \quad $pgp + adp \rightleftharpoons 3pg + atp$

\begin{equation}
    \kappa_7 = \frac{([pgp]/K_{S,pgp}^7)([adp]/K_{S,adp}^7)}%
                    {1 + ([pgp]/K_{S,pgp}^7)([adp]/K_{S,adp}^7) + ([3pg]/K_{P,3pg}^7)([atp]/K_{P,atp}^7)}
    \qquad
    \gamma_7 = 1 - \frac{[3pg]\,[atp]}{[pgp]\,[adp]\,K'_{eq,7}}
\end{equation}
\begin{equation}
    v_7 = Pgk \cdot k^+_{cat,7} \cdot \kappa_7 \cdot \gamma_7
    \label{eqn: pgk}
\end{equation}

% ---------------------------------------------------------------
\textbf{Phosphoglycerate mutase (GPM)} \quad EC: 5.4.2.11 \quad $3pg \rightleftharpoons 2pg$

\begin{equation}
    \kappa_8 = \frac{[3pg]/K_{S,3pg}^8}%
                    {1 + [3pg]/K_{S,3pg}^8 + [2pg]/K_{P,2pg}^8}
    \qquad
    \gamma_8 = 1 - \frac{[2pg]}{[3pg]\,K'_{eq,8}}
\end{equation}
\begin{equation}
    v_8 = GpmA \cdot k^+_{cat,8} \cdot \kappa_8 \cdot \gamma_8
    \label{eqn: gpm}
\end{equation}

% ---------------------------------------------------------------
\textbf{Enolase (ENO)} \quad EC: 4.2.1.11 \quad $2pg \rightleftharpoons pep$

\begin{equation}
    \kappa_9 = \frac{[2pg]/K_{S,2pg}^9}%
                    {1 + [2pg]/K_{S,2pg}^9 + [pep]/K_{P,pep}^9}
    \qquad
    \gamma_9 = 1 - \frac{[pep]}{[2pg]\,K'_{eq,9}}
\end{equation}
\begin{equation}
    v_9 = Eno \cdot k^+_{cat,9} \cdot \kappa_9 \cdot \gamma_9
    \label{eqn: eno}
\end{equation}
The stoichiometric matrix of this system is:
%\usepackage{blkarray}   % in preamble

\begin{align}
\mathbf{S} =
\begin{blockarray}{c ccccccccc}
       & \text{pts} & \text{pgi} & \text{pfk} & \text{fba} & \text{tpi} & \text{gap} & \text{pgk} & \text{gpm} & \text{eno} \\
\begin{block}{c[ccccccccc]}
  \text{2pg}  &  0 &  0 &  0 &  0 &  0 &  0 &  0 &  1 & -1 \\
  \text{3pg}  &  0 &  0 &  0 &  0 &  0 &  0 &  1 & -1 &  0 \\
  \text{dhap} &  0 &  0 &  0 &  1 & -1 &  0 &  0 &  0 &  0 \\
  \text{f6p}  &  0 &  1 & -1 &  0 &  0 &  0 &  0 &  0 &  0 \\
  \text{fbp}  &  0 &  0 &  1 & -1 &  0 &  0 &  0 &  0 &  0 \\
  \text{g3p}  &  0 &  0 &  0 &  1 &  1 & -1 &  0 &  0 &  0 \\
  \text{g6p}  &  1 & -1 &  0 &  0 &  0 &  0 &  0 &  0 &  0 \\
  \text{pep}  & -1 &  0 &  0 &  0 &  0 &  0 &  0 &  0 &  1 \\
  \text{pgp}  &  0 &  0 &  0 &  0 &  0 &  1 & -1 &  0 &  0 \\
\end{block}
\end{blockarray}
\label{eqn: S matrix}
\end{align}
\newpage

\section{Figures}
\subsection{\textit{Escherichia coli} metabolism}
\begin{figure}[h!]
    \centering
    \includegraphics[width=0.85\linewidth]{\SIdir figures/glycolisis_modified.png}
    \caption{Carbon metabolism pathways of \textit{E.coli} \cite{Ishii2007a}. The reaction highlighted in green encompass the studied reactions.}
    \label{fig:e.glicolisis}
\end{figure}
\clearpage

\subsection{Parameter selection}
\begin{figure}[h!]
    \centering
    \includegraphics[width=\linewidth]{\SIdir figures/figS_selection_path.pdf}
    \caption{\textbf{Parameter selection, decision by decision.} \textbf{a)} Parameter fixed by each decision (names coloured by type: PTS, $k_{cat}$, $K$) and the rule that fixed it: \SInStructuralZeros{} structural zeros of $\mathbf{R}$ in the structural phase, then \SInRankFixes{} directions of a rank-deficient FIM $\mathbf{F}$ and \SInCorrelationFixes{} highly correlated pair ($|r|>\SIcorrelation$) in the practical phase, one decision per iteration. The dotted lines separate the two phases and, lighter, the stages of the practical phase named above the panel: first the rank deficiency of $\mathbf{F}$ (the smallest singular value of $\mathbf{G}$ below the threshold), then, once $\mathbf{F}$ has full rank, the correlation criterion. \textbf{b)} Number of free parameters after each decision. \textbf{c)} Objective $g$ of the current fit after each decision; the structural fixes are refitted once, after the last one. Shaded: objective of the final model plus the 95\% likelihood-ratio threshold $\Delta=\SIthreshold$. \textbf{d)} Smallest singular value ratio $\sigma_{\min}/\sigma_1$ of $\mathbf{R}$ before each structural decision (threshold $\epsilon_1=\SIepsOne$; the open marker is the check that ends the structural phase) and smallest eigenvalue ratio $\lambda_{\min}/\lambda_{\max}$ of $\mathbf{F}$ at each practical iteration (threshold $\epsilon_2=\SIepsTwo$). \textbf{e)} Largest absolute correlation between free parameters once $\mathbf{F}$ has full rank.}
    \label{fig:selection_path}
\end{figure}

\begin{figure}[h!]
    \centering
    \includegraphics[width=\linewidth]{\SIdir figures/figS_final_identifiability.pdf}
    \caption{\textbf{Identifiability of the final model.} \textbf{a)} Singular values of the stacked rate sensitivity matrix $\mathbf{R}$ \cite{stigter2015fast} at the initial estimation (\SInAll{} free parameters) and at the final model (\SInFree{} free parameters); the \SInStructuralZeros{} values of the initial estimation below $\epsilon_1$ are its structural zeros. \textbf{b)} Eigenvalues of the FIM $\mathbf{F}$ after the structural phase (\SInStructural{} free parameters) and at the final model; the values below $\epsilon_2$ are the directions removed in the practical phase, and all eigenvalues of the final model lie above it (full rank). \textbf{c)} Correlations of the free parameters of the final model from $\mathbf{F}^{-1}$, under a header with the reaction of each parameter (colours of Figure~\ref{fig:main}a) and with names coloured by type; values are shown where $|r|\geq 0.5$, and none exceeds $\SIcorrelation$.}
    \label{fig:final_identifiability}
\end{figure}
\clearpage

\subsection{Model evaluation}
\begin{figure}[h!]
    \centering
    \includegraphics[width=\linewidth]{\SIdir figures/figS_parity_stages.pdf}
    \caption{\textbf{Parity of the model predictions at every stage.} Measured against predicted LI fluxes (top) and metabolite concentrations (bottom, log scale) for the literature model (all \SInAll{} parameters at their literature values), the initial estimation (all \SInAll{} parameters free), the final model (\SInFree{} free parameters) and LOOCV (held-out predictions of the final structure). RMSE and $R^2$ against the per-output mean of the measurements.}
    \label{fig:paritylitfinal}
\end{figure}

\begin{figure}[h!]
    \centering
    \includegraphics[width=\linewidth]{\SIdir figures/figS_parameter_evolution.pdf}
    \caption{\textbf{Parameter values across the stages.} All \SInAll{} parameters, grouped by reaction (coloured bands and strips, as in Figure~\ref{fig:main}a; names coloured by type), in the literature model, after the initial estimation and in the final model; the grey line spans the three values. Filled diamonds: free in the final model; open diamonds: fixed, at the value estimated when the parameter was fixed. Right: free, or the rule that fixed the parameter (with the practical iteration), in the rule colours of Figure~\ref{fig:selection_path}. The values are listed in Table~\ref{tab:param_evolution}.}
    \label{fig:fixed evol}
\end{figure}
\clearpage

\subsection{Parameter uncertainty}
\begin{figure}[h!]
    \centering
    \includegraphics[width=\linewidth]{\SIdir figures/figS_bootstrap_all.pdf}
    \caption{\textbf{Bootstrap distributions of the free parameters.} Case bootstrap with \SIbootN{} resamples of the \SInConditions{} conditions (\SIbootConverged{} converged), each refitted \SIbootStarts{}. $\theta^*$: estimate of a resample; $\hat\theta$: estimate of the final model (solid line). Shaded: 95\% percentile interval; panel titles: ratio of its ends $p_{97.5}/p_{2.5}$, and ``bound'' where the interval reaches a parameter bound. Colour: verdict of the bootstrap alone, practically identifiable if the ratio is at most \SIbootRatio{} and the interval does not reach a bound (Table~\ref{tab:free_parameters}).}
    \label{fig:boostrapdist}
\end{figure}

\begin{figure}[h!]
    \centering
    \includegraphics[width=\linewidth]{\SIdir figures/figS_profiles_all.pdf}
    \caption{\textbf{Profile likelihoods of the free parameters.} Objective of the fit with the parameter fixed at $\theta$ and the other free parameters re-optimised (warm started at the previous point, plus 3 starts around it; a point that a point farther out improved on by more than $\Delta$ was refitted from that point), relative to the best objective found, from $\hat\theta/\SIprofileFactor$ to $\SIprofileFactor\,\hat\theta$ with \SIprofilePoints{} logarithmically spaced points on each side. Dashed: 95\% threshold $\Delta=\chi^2_{1,0.95}/(2N)=\SIthreshold$. Missing points are failed re-optimisations or values beyond a parameter bound. Title and colour: profile class, in the colours of Figure~\ref{fig:main}e (an arm crosses $\Delta$ where it exceeds it and stays above it farther out; identifiable: both arms cross; one-sided: one arm; flat: none).}
    \label{fig:proflike}
\end{figure}

\clearpage
\section{Tables}

\begin{table}[h]
\centering
\caption{Parameter descriptions}
\label{table:Parameter descriptions}
\begin{tabular}{l|l}
\textbf{Symbol} & \textbf{Meaning} \\
\hline
$e_r$                  & Enzyme concentration of reaction $r$ \\[4pt]
$k^+_{cat,r}$          & Forward turnover number of reaction $r$ \\[4pt]
$K^r_{S,s}$            & Michaelis-Menten constant of substrate $s$
 in reaction $r$ \\[4pt]
$K^r_{P,p}$            & Michaelis-Menten constant of product $p$
 in reaction $r$\\[4pt]
$K'_{eq,r}$            & Equilibrium constant of reaction $r$\\[4pt]
$[s]$, $[p]$           & Concentration of substrate $s$ and product $p$\\[4pt]
$v^{\max}_{PTS}$, $K_{a1}$, $K_{a2}$, $K_{a3}$, $K_{g6p}$ & Parameters of the PTS rate law \\[4pt]
$n_{g6p}$              & Hill coefficient of the G6P inhibition of the PTS (fixed at 4)
\end{tabular}
\end{table}




\begin{table}[h]
\centering
\caption{Metabolites involved, abbreviation and  their mass-balance classification}
\label{tab: metabolites}
\begin{tabular}{lll}
\textbf{Abbrev.} & \textbf{Metabolite} & \textbf{Classification} \\
\hline
$glc$  & Glucose                    & Imbalanced \\
$g6p$  & Glucose 6-phosphate        & Balanced \\
$f6p$  & Fructose 6-phosphate       & Balanced \\
$fbp$  & Fructose 1,6-bisphosphate  & Balanced \\
$dhap$ & Dihydroxyacetone phosphate & Balanced \\
$g3p$  & Glyceraldehyde 3-phosphate & Balanced \\
$pgp$  & 1,3-Bisphosphoglycerate    & Balanced \\
$3pg$  & 3-Phosphoglycerate         & Balanced \\
$2pg$  & 2-Phosphoglycerate         & Balanced \\
$pep$  & Phosphoenolpyruvate        & Balanced \\
$pyr$  & Pyruvate                   & Imbalanced \\
$atp$  & ATP                        & Imbalanced \\
$adp$  & ADP                        & Imbalanced \\
$nad$  & NAD$^+$                    & Imbalanced \\
$nadh$ & NADH                       & Imbalanced \\
$pi$   & Inorganic phosphate        & Imbalanced \\
\end{tabular}
\end{table}


% con esto se saca la matriz de covarianza luego la de correlacion y se hace un heat map, y a parte se hace una matriz de sensibilidad con el mismo G (derivadas de los metabolitos (x) con respecto a los parametros) que usamos para la FIM. Con esto evaluamos los parametros a fijar, al fijarlos (ojalas muchos) corremos el modelo reducido para estimar parametros. Obtenemos intervalos de confianza entre estimaciones (las da el mismo algoritmo) Y sacamos conclusiones a partir de los resultados, como identificabilidad, comportamiento de vía, importancia de la via para el metabolismo de carbono celular, definimos el future work, que es seguir validando el modelo (cross validation) y la capacidad predictiva y lo del control coefficients Jv Jx, para sugerir cambios en enzima para aumentar o disminuir flujos. (sensibilidad de flujos ante concentracion de cada enzima)


% equilibrator (ph 7 pmg 3 ionic strenght 0.1) explicar que se utilizan los datos encontrados com parametros a fijar y los no encontrados se optimizan, upper y lower bounds se dejan en un rango grande, k cat y Km a partir de literatura (graficos de paper profe (rangos de enzima)) 



\renewcommand{\arraystretch}{1.2}

\begin{longtable}{c
>{\RaggedRight\arraybackslash}p{4cm}
>{\RaggedRight\arraybackslash}p{5.5cm}
>{\RaggedRight\arraybackslash}p{5.6cm}
>{\RaggedRight\arraybackslash}p{2cm}}

\caption{Glycolysis kinetic parameters in  \textit{E. coli} obtained from EcoCyc\cite{karp2025ecocyc}, Kadir et al. \cite{Kadir2010ModelingsimulationmainmetabolismEscherichiacoliitsseveralsinglegeneknockoutmutantsexperimentalverification} and Equilibrator \cite{beber2022equilibrator}.}
\label{tab:glicolisis}\\

\toprule
\textbf{\#} &
\textbf{Enzime (gen)} &
\textbf{Reaction} &
\textbf{Parameters} &
\textbf{Source} \\
\midrule
\endfirsthead

\multicolumn{5}{c}{\tablename\ \thetable\ -- continuation} \\
\toprule
\textbf{\#} &
\textbf{Enzime (gen)} &
\textbf{Reaction} &
\textbf{Parameters} &
\textbf{Source} \\
\midrule
\endhead

\midrule
\multicolumn{5}{r}{Continues on next page}
\endfoot

\bottomrule
\endlastfoot

1 &
Phosphotransferase system &
$glc + pep \to g6p + pyr$ &\makecell[l]{
$v_{max}^{PTS}$: 25.739 mmol/(gDCW·h)\vspace{3pt}\\
$K_{a1}$: 1.0 mM\vspace{3pt}\\
$K_{a2}$: 0.01 mM\vspace{3pt}\\
$K_{a3}$: 1.0 mM\vspace{3pt}\\
$n_{g6p}$: 4\vspace{3pt}\\
$K_{g6p}$: 0.5 mM\vspace{3pt}\\}
&
\cite{Kadir2010ModelingsimulationmainmetabolismEscherichiacoliitsseveralsinglegeneknockoutmutantsexperimentalverification}
\\
\hline\vspace{3pt}

2 &
Phosphoglucose isomerase (\textit{Pgi}) &
$g6p\rightleftharpoons{f6p}$ &\makecell[l]{
$K_{eq}:0.381$\vspace{3pt}\\
$K_m^{g6p}$: 1.018 mM\vspace{3pt}\\
$K_m^{f6p}$: 0.078 mM\vspace{3pt}\\
$k_{cat}:--$
}
%hay otro valor de f6p (0.2) en el mismo Ecocyc pero este (0.078) es del mismo estudio que el de g6p.
&
\href{https://ecocyc.org/gene?orgid=ECOLI&id=PGLUCISOM}{EcoCyc: EG10702}
\\
\hline\vspace{3pt}
3 &
Phosphofructokinase (\textit{PfkB}) &
$f6p+atp\to fbp+adp$  &\makecell[l]{
$K_{eq}:341.298$\vspace{3pt}\\
$K_m^{f6p}$: 0.013 mM\vspace{3pt}\\
$K_m^{atp}$: 0.02 mM\vspace{3pt}\\
$K_m^{fbp}:$ 0.14 mM\vspace{3pt}\\
$K_m^{adp}: --$\vspace{3pt}\\
$k_{cat}:--$\vspace{3pt}\\}
&
\href{https://ecocyc.org/gene?orgid=ECOLI&id=EG10700}{EcoCyc: EG10700}
\\
\hline\vspace{3pt}
4 &
Aldolase (\textit{FbaA}) &
$fbp\rightleftharpoons g3p+dhap$ 
&\makecell[l]{
$K_{eq}:0.1616$\vspace{3pt}\\
$K_m^{fbp}$: 0.24 mM\vspace{3pt}\\
$K_m^{g3p}:--$\vspace{3pt}\\
$K_m^{dhap}:--$\vspace{3pt}\\
$k_{cat}:10.33 s^{-1} \to 0.95 \frac{mmol}{mg\cdot h}$\vspace{3pt}\\
}
%kcat se calcula con MW de secuencia: 39147 D
&
\href{https://ecocyc.org/gene?orgid=ECOLI&id=EG10282}{EcoCyc: EG10282}
\\ 
\hline 

5 &
TPI (\textit{TpiA}) &
$dhap\rightleftharpoons g3p$ 
&\makecell[l]{
$K_{eq}:0.113$\vspace{3pt}\\
$K_m^{dhap}$: 1.03 mM\vspace{3pt}\\
$K_m^{g3p}: --$\vspace{3pt}\\
$k_{cat}$: 9000 s$^{-1} \to 1201.2 \frac{mmol}{mg\cdot h}$\vspace{3pt} \\
} 
%para calcular kcat se ocupa el MW de secuencia 26972 D
&
\href{https://ecocyc.org/gene?orgid=ECOLI&id=EG11015}{EcoCyc: EG11015}
\\
\hline
6 &
GAPDH (\textit{GapA}) &
$g3p+nad+pi\rightleftharpoons pgp+nadh$ 
&\makecell[l]{
$K_{eq}:0.00026$\vspace{3pt}\\
$K_m^{g3p}$: 0.89 mM\vspace{3pt}\\
$K_m^{nad}$: 0.045 mM\vspace{3pt}\\
$K_m^{pi}$: 0.53 mM\vspace{3pt}\\
$K_m^{pgp}:--$\vspace{3pt}\\
$K_m^{nadh}:--$\vspace{3pt}\\
$k_{cat}:--$}
&
\href{https://ecocyc.org/gene?orgid=ECOLI&id=GAPDH-A-MONOMER}{EcoCyc: EG10367}
\\
\hline
7 &
Phosphoglycerate kinase (\textit{Pgk}) &
$pgp+adp\rightleftharpoons3pg+atp$ &
\makecell[l]{
$K_{eq}:1849.572$\vspace{3pt}\\
$K_m^{pgp}:--$\vspace{3pt}\\
$K_m^{adp}:--$\vspace{3pt}\\
$K_m^{3pg}:--$\vspace{3pt}\\
$K_m^{atp}:--$\vspace{3pt}\\
$k_{cat}:--$}
&
Equilibrator \cite{beber2022equilibrator}  
\\
\hline
8 &
Phosphoglycerate mutase (\textit{GpmA}) &
$3pg\rightleftharpoons2pg$ 
&\makecell[l]{
$K_{eq}:0.1874$\vspace{3pt}\\
$K_m^{3pg}$: 0.2 mM\vspace{3pt}\\
$K_m^{2pg}$: 0.19 mM\vspace{3pt}\\
% en el informe 2 k_cat=220 1/s
$k_{cat}$: 330 s$^{-1}$ $\to$ 41.602 $\frac{mmol}{mg\cdot h}$\vspace{3pt} \\}
%k cat se calculó usando el MW del polipéptido obtenido de secuencia: 28.56kD
&
\href{https://ecocyc.org/gene?orgid=ECOLI&id=EG11699}{EcoCyc: EG11699 }
\\
\hline
9 &
Enolase (\textit{Eno}) &
$2pg\rightleftharpoons pep$ 
&\makecell[l]{
$K_{eq}:4.5273$\vspace{3pt} \\
$K_m^{2PG}$: 0.1 mM\vspace{3pt} \\
$K_m^{pep}:--$\vspace{3pt} \\
$k_{cat}:--$
}
&
\href{https://ecocyc.org/gene?orgid=ECOLI&id=ENOLASE-MONOMER#RXNS}{EcoCyc: EG10258}
\\

\end{longtable}

\input{\SIdir generated/table_parameter_values}

\input{\SIdir generated/table_free_parameters}
